\documentclass{ximera}
\title{Partial Derivatives and Level Sets}

\newcommand{\pskip}{\vskip 0.1 in}

\begin{document}
\begin{abstract}
Introduction to partial derivatives and level sets.
\end{abstract}
\maketitle



\section{The Hanging Weight}
\begin{question} \label{ExMD3fetet39}

A block $B$ with weight $w = 5$ pounds hangs suspended from two ropes attached to the ceiling at points $A_1$, $A_2$ and making respective angles $\theta_1$ and $\theta_2$ with the vertical as shown below. 


\begin{onlineOnly}
    \begin{center}
\desmos{0fudotupws}{900}{600}
\end{center}
\end{onlineOnly}

\href{https://www.desmos.com/calculator/0fudotupws}{163: Hanging Weight and Tension 3}


\item Our next task is to find an expression for the magnitude $T_1 = |\overrightarrow{T_1}|$ of the tension in rope $BA_1$.

\begin{enumerate}
\item Use triangle trigonometry to find an expression for the function
\[
   T_1 = f(\theta_1, \theta_2) \, , \, 0 <  \theta_1 \leq \pi/2 \, , \,  0 <  \theta_2 \leq \pi/2.  
\]
that expresses the magnitude $T_1 = |\overrightarrow{T_1}|$ of the tension in rope $BA_1$ in terms of the angles $\theta_1$ and $\theta_2$. %Your expression should also include $w$ (the weight of the block in pounds), but for the remainder of the question we'll regard the weight as constant.

%Work in general, \emph{not} with the specific values of the parameters shown above.

\item Check your expression with the worksheet above by substituting $\theta_1 = 0.5$, and $\theta_2=0.6$.

\item Drag sliders $\theta_1$ above to see how the tension in rope $BA_1$ changes as $\theta_1$ changes. Summarize your findings.

\item Drag $\theta_1$ back to $\theta_1=0.5$. Then drag slider $\theta_2$ to see how the tension  in rope $BA_1$ changes as $\theta_2$ changes. Summarize your findings.

\item Some level sets of the function $f$ are shown below. The sets are spaced at intervals of $\Delta T_1 = 0.5$ pounds.


\begin{onlineOnly}
    \begin{center}
\desmos{iahtoxdhzu}{900}{600}
\end{center}
\end{onlineOnly}

\href{https://www.desmos.com/calculator/iahtoxdhzu}{163: Hanging Weight Level Curves}

\item Use the level sets to approximate the partial derivatives
\[
    \frac{\partial T_1}{\partial \theta_1} \Big|_{(\theta_1, \theta_2)=(0.5, 0.6)}
\]
and
\[
    \frac{\partial T_1}{\partial \theta_2} \Big|_{(\theta_1, \theta_2)=(0.5, 0.6)} .
\]
Include units.


\item Use your expression for the function $T_1 = f(\theta_1, \theta_2)$ to compute these partials.

\item Use the results of part (g) to approximate the change in $T_1$ as $\theta_1$ increases from $0.5$ radians to $0.55$ radians and $\theta_2$ is held constant at $0.6$ radians. Then compare this approximation with the actual change.

\item Use the results of part (g) to approximate the change in $T_1$ as $\theta_2$ decreases from $0.6$ radians to $0.55$ radians and $\theta_1$ is held constant at $0.5$ radians.. Then compare this approximation with the actual change.

%\item Use part (g) to approximate the 

%\item Use ideas of implicit differentiation to approximate the change in 

\end{enumerate}
\end{question}

\section{Distances on a Sphere}
The function
\[
 s = f(\theta,\phi) = \arccos(\sin\theta \cos \phi) \, , \, 0\leq \theta \leq 2\pi, -\pi/2\leq \phi \leq \pi/2 ,
\]
expresses the distance (in miles) from $(0,1,0)$ of a point $P$ on a sphere with radius $1$ miles in terms of its latitude $\theta$ (in radians) and longitude $\phi$ (in radians). The distance is measured along the sphere.

\begin{onlineOnly}
    \begin{center}
\desmosThreeD{diqbhi8jch}{900}{600}
\end{center}
\end{onlineOnly}

\href{https://www.desmos.com/3d/diqbhi8jch}{163: Distances on Sphere 3D}

Some level sets of the function $f$ are graphed below.

\begin{onlineOnly}
    \begin{center}
\desmos{njngjntwn0}{900}{600}
\end{center}
\end{onlineOnly}

\href{https://www.desmos.com/calculator/njngjntwn0}{163: Level Sets Distances on Sphere}


\begin{enumerate}
\item Drag point $P$ above to get a sense of how the distance changes as the latitude and longitude of $P$ change.

\item Use the level curves to approximate the partials
\[
    \frac{\partial s}{\partial \theta} \Big|_{(\theta, \phi)=(2,1)}
\]
and
\[
     \frac{\partial s}{\partial \phi} \Big|_{(\theta, \phi)=(2,1)} .
\]
Include units.

\item Use the expression for $f$ above to compute the exact values of these partials.

\item Use the results of (b) and (c) to approximate the latitude of the point on the sphere near $P$ with longitude $\theta=2.1$ radians that is the same distance from $(0,1,0)$ as $P$.

\item Use ideas of Math 151 and implicit differentiation to answer part (d).



\end{enumerate}

\end{document}